% Box in stile terminale.
% Riproduce il terminale in uso: foot con il tema Omarchy attivo
% (colori in ~/.local/state/omarchy/current/theme/foot.ini).
% Richiede: tcolorbox, xcolor, listings (gia' caricati nel preambolo).
%
% NOTA: il terminale usa JetBrains Mono, che funziona solo con
% XeLaTeX/LuaLaTeX. Il documento usa pdfLaTeX, quindi si usa Fira Mono
% (stessa larghezza dei caratteri, forme simili).
%
% Uso:
%   \begin{terminalbox}
%   (*@\tprompt[main]{domain}@*)git status
%   On branch main
%   (*@\textcolor{termred}{errore in rosso}@*)
%   (*@\tprompt[main]{domain}@*)(*@\tcursor@*)
%   \end{terminalbox}
% Il testo e' verbatim. Per i colori si scrive LaTeX tra (*@ e @*).

\tcbuselibrary{skins,breakable,listings}

% T1 serve a Fira Mono per i simboli < > | _. OT1 resta l'encoding
% di default perche' e' l'ultima opzione.
\usepackage[T1,OT1]{fontenc}

% Fira Mono cambia \ttdefault: si salva il valore originale (Courier)
% e lo si ripristina, cosi' \texttt nel resto del documento non cambia.
\let\termoldttdefault\ttdefault
\usepackage[lining]{FiraMono}
\let\ttdefault\termoldttdefault

% Colori del tema (foot.ini)
\definecolor{termbg}{HTML}{111C18}
\definecolor{termfg}{HTML}{C1C497}
\definecolor{termbright}{HTML}{F7E8B2}
\definecolor{termaccent}{HTML}{509475}
\definecolor{termcyan}{HTML}{2DD5B7}
\definecolor{termmuted}{HTML}{53685B}
\definecolor{termred}{HTML}{FF5345}
\definecolor{termgreen}{HTML}{63B07A}
\definecolor{termyellow}{HTML}{E5C736}

% Font del terminale: 8 pt con interlinea 10.5 pt, circa 78 colonne
% nella larghezza del testo.
\newcommand{\termfont}{%
  \fontencoding{T1}\fontfamily{FiraMono-TLF}\fontsize{8}{10.5}\selectfont}

\newtcblisting{terminalbox}[1][]{%
  enhanced,
  breakable,
  listing only,
  sharp corners,
  colback=termbg,
  colframe=termaccent,
  boxrule=1.5pt,
  boxsep=0pt,
  left=10pt, right=10pt, top=10pt, bottom=5pt,
  listing options={%
    basicstyle=\termfont\color{termfg},
    columns=fixed,
    keepspaces=true,
    showstringspaces=false,
    breaklines=true,
    breakindent=0pt,
    breakautoindent=false,
    upquote=true,
    escapeinside={(*@}{@*)},
    % listings non legge l'UTF-8: ogni carattere non ASCII usato nei
    % box va elencato qui.
    extendedchars=true,
    literate=%
      {à}{{\`a}}1 {è}{{\`e}}1 {é}{{\'e}}1 {ì}{{\`i}}1 {ò}{{\`o}}1 {ù}{{\`u}}1
      {À}{{\`A}}1 {È}{{\`E}}1 {É}{{\'E}}1 {Ì}{{\`I}}1 {Ò}{{\`O}}1 {Ù}{{\`U}}1
      {’}{{\textquoteright}}1 {‘}{{\textquoteleft}}1
      {“}{{\textquotedblleft}}1 {”}{{\textquotedblright}}1,
  },
  #1
}

% Prompt come starship: directory, branch (facoltativo) e simbolo.
% Uso: \tprompt{domain} oppure \tprompt[main]{domain}
\newcommand{\tprompt}[2][]{%
  \textcolor{termcyan}{\textbf{#2}}%
  \if\relax\detokenize{#1}\relax\else\ \textcolor{termcyan}{\textit{#1}}\fi
  \ \textcolor{termcyan}{\textbf{>}}\ }

% Cursore a blocco, largo un carattere e alto una riga.
\newcommand{\tcursor}{%
  {\color{termbright}\rule[-0.3\baselineskip]{0.6em}{\baselineskip}}}

% Testo spento (come l'output "dim" del terminale).
\newcommand{\tdim}[1]{\textcolor{termmuted}{#1}}
